\documentclass[11pt]{article}
\usepackage{fontspec}
\usepackage[french]{babel}
\usepackage{amsmath}
\usepackage{amssymb}
\usepackage{amsthm}
\usepackage{mathrsfs}
\usepackage{tikz-cd}
\usepackage{enumitem}

\newcommand{\printedpage}[1]{\marginpar{\footnotesize p.~#1}}
\newcommand{\msp}[1]{\marginpar{\footnotesize [#1]}}
\newcommand{\Asph}{\mathrm{Asph}}
\newcommand{\Cof}{\mathrm{Cof}}

\begin{document}

\printedpage{175}
% head: A. Grothendieck : Dérivateurs, Ch. XII, Édition des Universités de Montpellier II et Paris VII
% foot: Version 30 mars 2003

% dangling arrows (from X_2, X_3, Y_2, Y_3, Z_2, Z_3) as printed
\[
\begin{tikzcd}[ampersand replacement=\&, column sep=small, row sep=small]
 \& \& {} \& \& {} \& \& U \arrow[dd, "\mu"] \\
 \& X_2 \arrow[rr] \arrow[dd] \arrow[ur, shorten >=1ex] \& \& X_3 \arrow[rr] \arrow[dd] \arrow[ur, shorten >=1ex] \& \& T \arrow[ur] \arrow[dd] \& \\
X_0 \arrow[rr, crossing over] \arrow[ur] \arrow[dd] \& \& X_1 \arrow[rr, crossing over] \arrow[ur] \arrow[dd, crossing over] \& \& S \arrow[ur, "\lambda"] \arrow[dd, crossing over] \& \& U' \arrow[dd] \\
 \& Y_2 \arrow[rr] \arrow[dd] \arrow[ur, shorten >=1.5ex] \& \& Y_3 \arrow[rr] \arrow[dd] \arrow[ur, shorten >=1.5ex] \& \& T' \arrow[ur] \arrow[dd] \& \\
Y_0 \arrow[rr, crossing over] \arrow[ur] \arrow[dd] \& \& Y_1 \arrow[rr, crossing over] \arrow[ur] \arrow[dd, crossing over] \& \& S' \arrow[ur, "\lambda'"] \arrow[dd, crossing over] \& \& U'' \\
 \& Z_2 \arrow[rr] \arrow[ur, shorten >=1.5ex] \& \& Z_3 \arrow[rr] \arrow[ur, shorten >=1.5ex] \& \& T'' \arrow[ur] \& \\
Z_0 \arrow[rr] \arrow[ur] \& \& Z_1 \arrow[rr] \arrow[ur] \& \& S'' \arrow[ur, "\lambda''"] \& \&
\end{tikzcd}
\]

La condition que $X$ soit $W$-cocartésien s'exprime par $\lambda \in W$, ou encore par $U \in \Asph_W$ (car $U$ s'identifie à la cofibre homotopique $Z(\lambda)$ de $\lambda$ \dots\,), de même $Y$ $W$-cocartésien s'exprime par $U' \in \Asph_W$, et $Z$ $W$-cocartésien par $U'' \in \Asph_W$. Or $U'' \sim Z(\mu)$, où $\mu : U \longrightarrow U'$ est dans $W$ si $U$, $U'$ sont $W$-asphériques, donc $U'' = Z(\mu)$ est alors $W$-asphérique. Cela prouve le lemme.

Cela montre que, moyennant W(7 bis) (en fait, on a utilisé seulement la partie ‘universelle’ de W(7 bis)), savoir que dans
\begin{tikzcd}[ampersand replacement=\&, column sep=small, row sep=small, baseline=(\tikzcdmatrixname-2-1.base)]
T_0 \arrow[rr] \arrow[dr, "i"'] \& \& T_2 \arrow[dr, "i'"] \& \\
 \& T_1 \arrow[rr] \& \& T_3
\end{tikzcd}
diagramme cocartésien d'immersions ouvertes, $i \in W \Longrightarrow i' \in W$ (au lieu de $\Longleftrightarrow$),

\msp{page 270}

la validité de W(7) se ramène au cas où le carré $Y$ est le carré trivial
\begin{tikzcd}[ampersand replacement=\&, column sep=small, row sep=small, baseline=(\tikzcdmatrixname-2-1.base)]
 \& e \arrow[dr] \& \\
e \arrow[ur] \arrow[dr] \& \& e, \\
 \& e \arrow[ur] \&
\end{tikzcd}
$e$ l'objet final de Cat.

On peut maintenant conclure : On a (moyennant la partie directe de W(7 bis))
\[
\mathrm{a}' \ \& \ \mathrm{b}' \Longrightarrow \mathrm{c}',
\]
car si $Z_0$, $Z_1$ sont $\in \Asph_W$, alors $i_1 \in W$, donc $j_1 \in W$, et si de plus $Z_2 \in \Asph_W$, on a donc $Z_3 \in \Asph_W$.
\[
\begin{tikzcd}[ampersand replacement=\&, column sep=small, row sep=small]
 \& Z_1 \arrow[dr, "j_2"] \& \\
Z_0 \arrow[ur, "i_1"] \arrow[dr, "i_2"'] \& \& Z_3 \\
 \& Z_2 \arrow[ur, "j_1"'] \&
\end{tikzcd}
\]

Si on a aussi la partie inverse (‘inhabituelle’) de W(7 bis), alors
\[
\mathrm{b}' \ \& \ \mathrm{c}' \Longrightarrow \mathrm{a}',
\]
car si $Z_1$ et $Z_3$ sont $W$-asphériques, alors $j_2 \in W$, donc $i_2 \in W$, et si de plus $Z_2$ est $W$-asphérique, on conclut $Z_0$ $W$-asphérique.

\printedpage{176}
% head: A. Grothendieck : Dérivateurs, Ch. XII, Édition des Universités de Montpellier II et Paris VII
% foot: Version 30 mars 2003

Enfin, supposons que $W$ soit \textit{stable par facteurs directs,} et W(7 bis) satisfait sous forme directe. Je dis qu'alors
\[
\mathrm{c}' \ \& \ \mathrm{a}' \Longrightarrow \mathrm{b}',
\]
i.e.
\[
Z_0 \text{ et } Z_3 \ \in \Asph_W \qquad\quad \Longrightarrow \qquad\quad Z_1 \text{ et } Z_2 \ \in \Asph_W.
\]
Factorisant $i_1$ et $i_2$ en
\[
\begin{aligned}
Z_0 &\xrightarrow{i'_1} Z'_1 \xrightarrow{\alpha_1} Z_1, \\
Z_0 &\xrightarrow{i'_2} Z'_2 \xrightarrow{\alpha_2} Z_2,
\end{aligned}
\]
avec $i'_1$, $i'_2$ des immersions ouvertes de Dwyer, et $\alpha_1, \alpha_2 \in W$, on est ramené au cas où

\msp{page 271}

$i_1$, $i_2$ sont des \textit{immersions ouvertes de Dwyer,} et où le carré est cocartésien (ce qui implique qu'il est $W$-cocartésien). Faisant alors le cochangement de base par $Z_0 \xrightarrow{p_0} Z'_0 = e$, on trouve un homomorphisme de carrés cocartésiens
\[
\begin{tikzcd}[ampersand replacement=\&, column sep=small, row sep=small]
 \& Z_1 \arrow[rr] \arrow[dd, "p_1" near end] \& \& Z_3 \arrow[dd, "p_3"] \\
Z_0 \arrow[ur] \arrow[rr, crossing over] \arrow[dd, "p_0"'] \& \& Z_2 \arrow[ur] \& \\
 \& Z'_1 \arrow[rr] \& \& Z'_3, \\
e = Z'_0 \arrow[ur] \arrow[rr] \& \& Z'_2 \arrow[ur] \arrow[from=uu, crossing over, "p_2"' near end] \&
\end{tikzcd}
\]
où toutes les faces du cube sont $W$-cocartésiennes, et comme $p_0 \in W$ ($Z_0$ étant $W$-asphérique), les autres flèches verticales $p_1$, $p_2$, $p_3$ sont \texttt{[dans]} $W$, donc le carré $Z'_*$ satisfait aussi $Z'_0$ et $Z'_3$ \texttt{[}$W$\texttt{-]}asphériques (et de plus $Z'_0 = e$), et il faut prouver $Z'_1$, $Z'_2$ $W$-asphériques (d'où itou pour $Z_1$, $Z_2$). Donc, on est ramené au cas particulier où $Z_0 = e$.

Mais dans ce cas il existe une rétraction (triviale!) de $Z_1$ sur $Z_0$, donc aussi de $Z_3$ sur $Z_2$. Ainsi, $Z_2$ est un facteur

\msp{page 272}

direct d'un objet $W$-asphérique, ce qui implique, par l'hypothèse sur $W$, que $Z_2$ est $W$-asphérique. Par symétrie, $Z_1$ l'est aussi. On a gagné!

\textbf{NB} En fait, ces arguments sont généraux, et donnent ceci :

\setcounter{footnote}{136}
\textbf{Théorème.} \textit{Soit $(\mathcal{M}, W)$ une catégorie des modèles, admettant une théorie de la cofibre $W$-homotopique \textup{(}\footnote{Mais on ne sait vérifier l'existence d'une telle théorie dans $(\mathrm{Cat}, W)$ que si on admet la forme faible de W(7), i.e. $f_0, f_1, f_2 \in W \implies f_3 \in W$, ou quelque cas particulier de celle-ci. Mais le \textit{lemme de Brown} sauve la situation, cf. XV \S3.}\textup{)} (i.e. on a une factorisation de toute flèche $f : X \longrightarrow Y$ en $pi$, avec $i \in \Cof_W$, $p \in W$.) On a donc une notion de carré $W$-cocartésien dans $\mathcal{M}$ \textup{(}\footnote{à vérifier que la factorisation suffit, même si tous les objets ne sont pas fibrants sur $e$. Mais on peut alors remplacer $\mathcal{M}$ par $\mathcal{M}_f$, et développer ce formalisme des fibres homotopiques en conséquence \dots}\textup{)}. Soit
\printedpage{177}
% head: A. Grothendieck : Dérivateurs, Ch. XII, Édition des Universités de Montpellier II et Paris VII
% foot: Version 30 mars 2003
alors $f : X \longrightarrow Y$ un homomorphisme d'un carré $W$-cocartésien $X$ dans un autre $Y$, $f = (f_0, f_1, f_2, f_3)$.}
\[
\begin{tikzcd}[ampersand replacement=\&, column sep=small, row sep=small]
 \& X_2 \arrow[rr] \arrow[dd, "f_2" near end] \& \& X_3 \arrow[dd, "f_3"] \\
X_0 \arrow[ur] \arrow[rr, crossing over] \arrow[dd, "f_0"'] \& \& X_1 \arrow[ur] \& \\
 \& Y_2 \arrow[rr] \& \& Y_3 \\
Y_0 \arrow[ur] \arrow[rr] \& \& Y_1 \arrow[ur] \arrow[from=uu, crossing over, "f_1"' near end] \&
\end{tikzcd}
\]

\textit{Considérons les conditions sur $f$}

\begin{enumerate}
\item[\textit{a)}] \textit{$f_0 \in W$.}
\item[\textit{b)}] \textit{$f_1$ et $f_2 \in W$.}
\item[\textit{c)}] \textit{$f_3 \in W$.}
\end{enumerate}

\textit{Alors}

\begin{enumerate}
\item[\textit{1°)}] \textit{a \& b $\Longrightarrow$ c.}

\msp{page 273}

\item[\textit{2°)}] \textit{Supposons que $\mathcal{M}$ ait un objet final $e$. Alors l'implication}
\[
\textit{b \& c} \Longrightarrow \textit{a}
\]
\textit{pour}~tout \textit{homomorphisme $f : X \longrightarrow Y$ de carrés $W$-cocartésiens équivaut à l'axiome suivant sur $W$ :}

Axiome des carrés $W$-cocartésiens conservatifs (\footnote{cf. argument page 78 qui montre que l'implication b \& c $\Longrightarrow$ a entraine l'axiome conservatif.}) : \textit{Si}
\[
\begin{tikzcd}[ampersand replacement=\&, column sep=small, row sep=small]
 \& Z_1 \arrow[dr, "j_2"] \& \\
Z_0 \arrow[ur, "i_1"] \arrow[dr, "i_2"'] \& \& Z_3 \\
 \& Z_2 \arrow[ur, "j_1"'] \&
\end{tikzcd}
\]
\textit{est $W$-cocartésien, alors $j_1 \in W \Longrightarrow i_1 \in W$ (donc en fait $i_1 \in W \Longleftrightarrow j_1 \in W$ et $i_2 \in W \Longleftrightarrow j_2 \in W$).}

\textit{Cet axiome est équivalent à la conjonction des deux axiomes suivants.}

\textbf{Cons 1.} \textit{Si on a un carré cocartésien}
\printedpage{178}
% head: A. Grothendieck : Dérivateurs, Ch. XII, Édition des Universités de Montpellier II et Paris VII
% foot: Version 30 mars 2003
\[
\begin{tikzcd}[ampersand replacement=\&, column sep=small, row sep=small]
S \arrow[rr] \arrow[dr] \& \& T \arrow[dr] \& \\
 \& e \arrow[rr] \& \& U,
\end{tikzcd}
\]
\textit{avec $f \in \Cof_W$, alors $U \in \Asph_W \Longrightarrow f \in W$ (donc $U \in \Asph_W \Longleftrightarrow f \in W$).}

\textbf{Cons 2.} \textit{Si on a un carré cartésien,}
\[
\begin{tikzcd}[ampersand replacement=\&, column sep=small, row sep=small]
 \& Z_1 \arrow[dr] \& \\
Z_0 \arrow[ur, "i_1"] \arrow[dr, "i_2"'] \& \& Z_3, \\
 \& Z_2 \arrow[ur] \&
\end{tikzcd}
\]
\textit{avec $i_1, i_2 \in \Cof_W$, alors : Si $Z_1$, $Z_2$, $Z_3$ sont $W$-asphériques, $Z_0$ l'est aussi.}

\msp{page 274}

\item[\textit{3°)}] \textit{Supposons que tout facteur direct d'un objet $W$-asphérique soit $W$-asphérique, ce qui est le cas notamment si $W$ est lui-même stable par facteurs directs, a fortiori si $W$ est fortement saturé. Supposons de plus que $\mathcal{M}$ ait} \texttt{[un]} \textit{objet final. Alors on a (sous les conditions générales du début)}
\[
\textit{c \& a} \Longrightarrow \textit{b,}
\]
\textit{i.e. si $f_0$, $f_3$ sont} \texttt{[dans]} \textit{$W$, de même $f_1$, $f_2$.}
\end{enumerate}

Cette dernière partie peut se déployer ainsi :

\textbf{Corollaire 1.} \textit{Sous l'hypothèse supplémentaire} \textup{3\,°)} \textit{sur $W$, la validité de c \& a $\Longrightarrow$ b (pour toute flèche $f$ entre carrés $W$-cocartésiens) équivaut à chacun des deux autres axiomes suivants sur $W$ (qui sont donc équivalents) :}

\begin{enumerate}
\item[\textit{1)}] \textit{Si on a un carré $W$-cocartésien}
\[
\begin{tikzcd}[ampersand replacement=\&, column sep=small, row sep=small]
 \& Z_1 \arrow[dr] \& \\
Z_0 \arrow[ur] \arrow[dr] \& \& Z_3, \\
 \& Z_2 \arrow[ur] \&
\end{tikzcd}
\]
\textit{alors $Z_0$ et $Z_3$ $W$-asphériques $\Longrightarrow$ $Z_1$ et $Z_2$ $W$-asphériques.}
\item[\textit{1$'$)}] \textit{Comme 1), avec $Z_0 = e$ (dans le cas où on suppose que $\mathcal{M}$ admet} \texttt{[un]} \textit{objet final $e$).}
\end{enumerate}

\msp{page 275}

\textbf{Corollaire 2.} \textit{Supposons que $W$ satisfasse l'axiome des carrés $W$-cocartésiens conservatifs, que tout facteur direct d'un objet $W$-asphérique soit $W$-asphérique (p. ex. $W$ stable par facteurs directs), enfin que $\mathcal{M}$ admette un objet final. Alors pout tout homomorphisme $f : X \longrightarrow Y$ de carrés $W$-cocartésiens, la conjonction de deux parmi les conditions a) b) c) de p. 272 entraîne la troisième.} % [sic] pout -> pour, typographical error

\end{document}
